\IfFileExists{stacks-project.cls}{%
\documentclass{stacks-project}
}{%
\documentclass{amsart}
}
\usepackage{amssymb}

% For dealing with references we use the comment environment
\usepackage{verbatim}
\newenvironment{reference}{\comment}{\endcomment}
%\newenvironment{reference}{}{}
\newenvironment{slogan}{\comment}{\endcomment}
\newenvironment{history}{\comment}{\endcomment}

% For commutative diagrams we use Xy-pic
\usepackage[all]{xy}

% We use 2cell for 2-commutative diagrams.
\xyoption{2cell}
\UseAllTwocells

% We use multicol for the list of chapters between chapters
\usepackage{multicol}

% This is generally recommended for better output
\usepackage{lmodern}
\usepackage[T1]{fontenc}
\usepackage{textalpha}

% For cross-file-references
\usepackage{xr-hyper}
\newcommand{\maybeexternaldocument}[2]{%
  \IfFileExists{#2.aux}{%
    \externaldocument[#1]{#2}%
  }{%
    \PackageWarning{stack}{Missing #2.aux; external references with prefix #1 unavailable}%
  }%
}
\let\externaldocumentorig\externaldocument
\renewcommand{\externaldocument}[2][]{%
  \IfFileExists{#2.aux}{%
    \externaldocumentorig[#1]{#2}%
  }{%
    \PackageWarning{stack}{Missing #2.aux; external references with prefix #1 unavailable}%
  }%
}

% Package for hypertext links:
\usepackage{hyperref}

% For any local file, say "hello.tex" you want to link to please
% use \externaldocument[hello-]{hello}
\externaldocument[introduction-]{introduction}
\externaldocument[conventions-]{conventions}
\externaldocument[sets-]{sets}
\externaldocument[categories-]{categories}
\externaldocument[topology-]{topology}
\externaldocument[sheaves-]{sheaves}
\externaldocument[sites-]{sites}
\externaldocument[stacks-]{stacks}
\externaldocument[fields-]{fields}
\externaldocument[algebra-]{algebra}
\externaldocument[brauer-]{brauer}
\externaldocument[homology-]{homology}
\externaldocument[derived-]{derived}
\externaldocument[simplicial-]{simplicial}
\externaldocument[more-algebra-]{more-algebra}
\externaldocument[smoothing-]{smoothing}
\externaldocument[modules-]{modules}
\externaldocument[sites-modules-]{sites-modules}
\externaldocument[injectives-]{injectives}
\externaldocument[cohomology-]{cohomology}
\externaldocument[sites-cohomology-]{sites-cohomology}
\externaldocument[dga-]{dga}
\externaldocument[dpa-]{dpa}
\externaldocument[sdga-]{sdga}
\externaldocument[hypercovering-]{hypercovering}
\externaldocument[schemes-]{schemes}
\externaldocument[constructions-]{constructions}
\externaldocument[properties-]{properties}
\externaldocument[morphisms-]{morphisms}
\externaldocument[coherent-]{coherent}
\externaldocument[divisors-]{divisors}
\externaldocument[limits-]{limits}
\externaldocument[varieties-]{varieties}
\externaldocument[topologies-]{topologies}
\externaldocument[descent-]{descent}
\externaldocument[perfect-]{perfect}
\externaldocument[more-morphisms-]{more-morphisms}
\externaldocument[flat-]{flat}
\externaldocument[groupoids-]{groupoids}
\externaldocument[more-groupoids-]{more-groupoids}
\externaldocument[etale-]{etale}
\externaldocument[chow-]{chow}
\externaldocument[intersection-]{intersection}
\externaldocument[pic-]{pic}
\externaldocument[weil-]{weil}
\externaldocument[adequate-]{adequate}
\externaldocument[dualizing-]{dualizing}
\externaldocument[duality-]{duality}
\externaldocument[discriminant-]{discriminant}
\externaldocument[derham-]{derham}
\externaldocument[local-cohomology-]{local-cohomology}
\externaldocument[algebraization-]{algebraization}
\externaldocument[curves-]{curves}
\externaldocument[resolve-]{resolve}
\externaldocument[models-]{models}
\externaldocument[functors-]{functors}
\externaldocument[equiv-]{equiv}
\externaldocument[pione-]{pione}
\externaldocument[etale-cohomology-]{etale-cohomology}
\externaldocument[proetale-]{proetale}
\externaldocument[relative-cycles-]{relative-cycles}
\externaldocument[more-etale-]{more-etale}
\externaldocument[trace-]{trace}
\externaldocument[crystalline-]{crystalline}
\externaldocument[spaces-]{spaces}
\externaldocument[spaces-properties-]{spaces-properties}
\externaldocument[spaces-morphisms-]{spaces-morphisms}
\externaldocument[decent-spaces-]{decent-spaces}
\externaldocument[spaces-cohomology-]{spaces-cohomology}
\externaldocument[spaces-limits-]{spaces-limits}
\externaldocument[spaces-divisors-]{spaces-divisors}
\externaldocument[spaces-over-fields-]{spaces-over-fields}
\externaldocument[spaces-topologies-]{spaces-topologies}
\externaldocument[spaces-descent-]{spaces-descent}
\externaldocument[spaces-perfect-]{spaces-perfect}
\externaldocument[spaces-more-morphisms-]{spaces-more-morphisms}
\externaldocument[spaces-flat-]{spaces-flat}
\externaldocument[spaces-groupoids-]{spaces-groupoids}
\externaldocument[spaces-more-groupoids-]{spaces-more-groupoids}
\externaldocument[bootstrap-]{bootstrap}
\externaldocument[spaces-pushouts-]{spaces-pushouts}
\externaldocument[spaces-chow-]{spaces-chow}
\externaldocument[groupoids-quotients-]{groupoids-quotients}
\externaldocument[spaces-more-cohomology-]{spaces-more-cohomology}
\externaldocument[spaces-simplicial-]{spaces-simplicial}
\externaldocument[spaces-duality-]{spaces-duality}
\externaldocument[formal-spaces-]{formal-spaces}
\externaldocument[restricted-]{restricted}
\externaldocument[spaces-resolve-]{spaces-resolve}
\externaldocument[formal-defos-]{formal-defos}
\externaldocument[defos-]{defos}
\externaldocument[cotangent-]{cotangent}
\externaldocument[examples-defos-]{examples-defos}
\externaldocument[algebraic-]{algebraic}
\externaldocument[examples-stacks-]{examples-stacks}
\externaldocument[stacks-sheaves-]{stacks-sheaves}
\externaldocument[criteria-]{criteria}
\externaldocument[artin-]{artin}
\externaldocument[quot-]{quot}
\externaldocument[stacks-properties-]{stacks-properties}
\externaldocument[stacks-morphisms-]{stacks-morphisms}
\externaldocument[stacks-limits-]{stacks-limits}
\externaldocument[stacks-cohomology-]{stacks-cohomology}
\externaldocument[stacks-perfect-]{stacks-perfect}
\externaldocument[stacks-introduction-]{stacks-introduction}
\externaldocument[stacks-more-morphisms-]{stacks-more-morphisms}
\externaldocument[stacks-geometry-]{stacks-geometry}
\externaldocument[moduli-]{moduli}
\externaldocument[moduli-curves-]{moduli-curves}
\externaldocument[examples-]{examples}
\externaldocument[exercises-]{exercises}
\externaldocument[guide-]{guide}
\externaldocument[desirables-]{desirables}
\externaldocument[coding-]{coding}
\externaldocument[obsolete-]{obsolete}
\externaldocument[fdl-]{fdl}
\externaldocument[index-]{index}

% Heap Project documents
\externaldocument[linear-algebra-]{linear-algebra}
\externaldocument[groups-]{groups}
\externaldocument[complex-analysis-]{complex-analysis}
\externaldocument[functional-analysis-]{functional-analysis}
\externaldocument[categories-]{categories}
\externaldocument[commutative-algebra-]{commutative-algebra}
\externaldocument[ringed-spaces-]{ringed-spaces}
\externaldocument[homological-algebra-]{homological-algebra}
\externaldocument[lie-algebras-]{lie-algebras}
\externaldocument[representation-theory-]{representation-theory}
\externaldocument[smooth-manifolds-]{smooth-manifolds}
\externaldocument[differential-topology-]{differential-topology}
\externaldocument[algebraic-topology-]{algebraic-topology}
\externaldocument[differential-geometry-]{differential-geometry}
\externaldocument[algebraic-geometry-]{algebraic-geometry}
\externaldocument[symplectic-geometry-]{symplectic-geometry}
\externaldocument[complex-geometry-]{complex-geometry}
\externaldocument[hodge-theory-]{hodge-theory}
\maybeexternaldocument{deformations-}{deformations}
\externaldocument[vertex-algebras-]{vertex-algebras}
\externaldocument[springer-]{springer}
\externaldocument[infty-categories-]{infty-categories}
\externaldocument[seminars-differential-geometry-]{%
seminars-differential-geometry}
\externaldocument[seminars-complex-geometry-]{%
seminars-complex-geometry}
\externaldocument[seminars-algebraic-geometry-]{%
seminars-algebraic-geometry}
\externaldocument[seminars-others-]{%
seminars-others}
%\externaldocument[geometric-prequantization-]{geometric-prequantization}
\externaldocument[surfaces-]{surfaces}
\externaldocument[birational-maps-commutative-algebra-]{birational-maps-commutative-algebra}
\externaldocument[cremona-transformations-]{cremona-transformations}
\externaldocument[derived-categories-of-sheaves-]{derived-categories-of-sheaves}
\externaldocument[geometric-invariant-theory-]{geometric-invariant-theory}
\externaldocument[geometric-stability-conditions-and-group-actions-]{geometric-stability-conditions-and-group-actions}
\externaldocument[moduli-spaces-of-sheaves-]{moduli-spaces-of-sheaves}
\externaldocument[bridgeland-stability-general-theory-]{bridgeland-stability-general-theory}
\externaldocument[hilbert-schemes-]{hilbert-schemes}
\externaldocument[noncommutative-abelian-surfaces-and-kummer-type-hyperkahler-varieties-]{noncommutative-abelian-surfaces-and-kummer-type-hyperkahler-varieties}
\externaldocument[mumford-tate-groups-in-hodge-theory-]{mumford-tate-groups-in-hodge-theory}
\externaldocument[hodge-birational-atoms-2026-]{hodge-birational-atoms-2026}
\externaldocument[xxii-egd-2026-]{%
xxii-egd-2026}
\externaldocument[pde-]{pde}
\externaldocument[probability-]{probability}

\externaldocument[linguistics-]{linguistics}
\externaldocument[genealogia-familiar-]{%
genealogia-familiar}

\externaldocument[%
16th-alga-meeting-2026-]{%
16th-alga-meeting-2026}

\externaldocument[%
iv-jornadas-lefschetz-2026-]{%
iv-jornadas-lefschetz-2026}

\externaldocument[reading-the-masters-]{%
reading-the-masters}

\externaldocument[real-analysis-]{%
real-analysis}

\externaldocument[optimal-transport-]{%
optimal-transport}

\externaldocument[history-]{history}

% Theorem environments.
%
\theoremstyle{plain}
\newtheorem{theorem}[subsection]{Theorem}
\newtheorem{proposition}[subsection]{Proposition}
\newtheorem{lemma}[subsection]{Lemma}

\theoremstyle{definition}
\newtheorem{definition}[subsection]{Definition}
\newtheorem{example}[subsection]{Example}
\newtheorem{exercise}[subsection]{Exercise}
\newtheorem{situation}[subsection]{Situation}

\theoremstyle{remark}
\newtheorem{remark}[subsection]{Remark}
\newtheorem{remarks}[subsection]{Remarks}

\numberwithin{equation}{subsection}

% Macros
%
\def\lim{\mathop{\mathrm{lim}}\nolimits}
\def\colim{\mathop{\mathrm{colim}}\nolimits}
\def\Spec{\mathop{\mathrm{Spec}}}
\def\Hom{\mathop{\mathrm{Hom}}\nolimits}
\def\Ext{\mathop{\mathrm{Ext}}\nolimits}
\def\SheafHom{\mathop{\mathcal{H}\!\mathit{om}}\nolimits}
\def\SheafExt{\mathop{\mathcal{E}\!\mathit{xt}}\nolimits}
\def\Sch{\mathit{Sch}}
\def\Mor{\mathop{\mathrm{Mor}}\nolimits}
\def\Ob{\mathop{\mathrm{Ob}}\nolimits}
\def\Sh{\mathop{\mathit{Sh}}\nolimits}
\def\NL{\mathop{N\!L}\nolimits}
\def\CH{\mathop{\mathrm{CH}}\nolimits}
\def\proetale{{pro\text{-}\acute{e}tale}}
\def\etale{{\acute{e}tale}}
\def\QCoh{\mathit{QCoh}}
\def\Ker{\mathop{\mathrm{Ker}}}
\def\Im{\mathop{\mathrm{Im}}}
\def\Coker{\mathop{\mathrm{Coker}}}
\def\Coim{\mathop{\mathrm{Coim}}}

% Boxtimes
%
\DeclareMathSymbol{\boxtimes}{\mathbin}{AMSa}{"02}

%
% Macros for moduli stacks/spaces
%
\def\QCohstack{\mathcal{QC}\!\mathit{oh}}
\def\Cohstack{\mathcal{C}\!\mathit{oh}}
\def\Spacesstack{\mathcal{S}\!\mathit{paces}}
\def\Quotfunctor{\mathrm{Quot}}
\def\Hilbfunctor{\mathrm{Hilb}}
\def\Curvesstack{\mathcal{C}\!\mathit{urves}}
\def\Polarizedstack{\mathcal{P}\!\mathit{olarized}}
\def\Complexesstack{\mathcal{C}\!\mathit{omplexes}}
% \Pic is the operator that assigns to X its picard group, usage \Pic(X)
% \Picardstack_{X/B} denotes the Picard stack of X over B
% \Picardfunctor_{X/B} denotes the Picard functor of X over B
\def\Pic{\mathop{\mathrm{Pic}}\nolimits}
\def\Picardstack{\mathcal{P}\!\mathit{ic}}
\def\Picardfunctor{\mathrm{Pic}}
\def\Deformationcategory{\mathcal{D}\!\mathit{ef}}


\begin{document}

\title{Salvador field notes}
\maketitle

\phantomsection
\label{section-salvador-fieldnotes}

\tableofcontents

\section{Purpose and source status}

These notes record the visit of
September 26, 2026.  They preserve
what was actually seen, the wording of
museum labels, the questions that arose,
and later explanations.

Four kinds of statements must remain
separate:

\begin{itemize}
\item direct observation or a photographed
label;
\item a museum or curatorial statement;
\item a later bibliographical verification;
\item a provisional interpretation or an
open question.
\end{itemize}

A museum label is evidence about what the
museum says.  It is not automatically the
last word on history, religion, language,
or geography.

\section{Route and negative evidence}

The places entered were the Cathedral
Basilica of Salvador, MUNCAB, and the
Casa do Benin.

The Church of Our Lady of the Rosary of
the Black People was found closed.  No
note should imply that its interior was
visited.

The Museum of Mercy was not visited that
day.  The visit was abandoned because it
was crowded.  Its material should not be
mixed into the record of the museums that
were actually seen.

\section{Cathedral Basilica}

A panel inside the Cathedral presented the
\emph{Caminho de Ant\^onio}, a six-kilometre
pilgrimage route between the Church of
Saint Anthony of Barra and the Sanctuary
of Saint Anthony Beyond the Carmo.

The same panel stated that the present
church was built from 1657 to 1672 as the
church of the Jesuit college and became a
cathedral in 1765.  These dates should be
read as the museum-panel account and
checked against architectural histories
before finer claims are made.

The visit prompted three questions:

\begin{itemize}
\item what distinguishes a cathedral from
an ordinary church;
\item who the Jesuits are;
\item how the colonial college relates to
modern Jesuit institutions such as
PUC--Rio and the Ibero in Mexico.
\end{itemize}

\section{MUNCAB and the Bienal}

The first photographed curatorial panel
introduced the Salvador itinerancy of the
36th Bienal de S\~ao Paulo,
\emph{Nem todo viandante anda estradas---
Da humanidade como pr\'atica}.

The panel described MUNCAB as a
``museum-territory''.  It asked how art
history produces presences, absences,
centralities, and silences.  It explicitly
rejected the idea that belonging requires
a homogeneous identity.

The intellectual constellation named by
the panel included Concei\c{c}\~ao Evaristo,
L\'elia Gonzalez, Beatriz Nascimento,
Milton Santos, bell hooks, and
\'Edouard Glissant.

This panel was not immediately persuasive
to Daniel.  The next exhibition, by Hariel
Revignet, generated a much more concrete
entry into the same questions.

\section{Hariel Revignet}

The exhibition was
\emph{Omi Tutu---\'Agua fresca}, by the
Brazilian-Gabonese artist Hariel Revignet,
with curatorial text by Jamile Coelho.

The sentence that most clearly called
attention was the image of a river entering
the sea.  The river becomes impossible to
distinguish, but its water has not ceased
to exist.  The curatorial text used this
image to think continuity through change,
not the intact survival of an unchanged
origin.

The text placed Oxum, Iemanj\'a, and Nan\~a
in a field of water, gestation, memory,
and renewal.  It described a feminine
principle that is not passive but creates,
sustains, receives, and transforms.

A photographed canvas was displayed beside
two labels:

\begin{itemize}
\item \emph{Ori Ej\'a}, 2025;
\item \emph{Iya Ej\'a}, 2025.
\end{itemize}

Both labels gave the medium as acrylic on
canvas with cowrie-shell embroidery.  The
photograph alone does not establish which
of the two adjacent labels belonged to the
particular canvas photographed.  The notes
must not silently assign it a title.

The canvas showed a dark female figure,
a rounded or pregnant abdomen, hands over
the body, a fish tail, and a dense surface
of cowries.  These are direct visual
observations.  Identifying the figure as a
particular orix\'a would require further
evidence.

The installation label read
\emph{Ziquid\'as---Rio de Mi\c{c}angas},
2026.  It listed strings of beads,
gold-coloured brass wire, a\c{c}a\'\i{}
seeds, small white clay vessels, tiger
cowries, a large clay vessel, gourds, water
jugs, a clay bowl, a basin, and seed-fish.

The hanging bead strings visually suggested
falling water.  The possible relation to
waist beads, the artist's grandmother, and
Gabonese family memory was an interpretive
lead raised in the conversation.  It should
be retained as a lead, not stated without a
source as the meaning of the installation.

\section{Diogum and the ofa}

A small iron sculpture was displayed beside
an artist biography headed ``Di Ogum''.
The text identified the artist as Diogo
Oliveira and mentioned symbols such as the
ofa of Ox\'ossi, the axe of Xang\^o, and the
trident of Exu.

Daniel correctly recognized that the
pictured curved form with a central shaft
looked more like the bow and arrow of
Ox\'ossi than a symbol of Ogum.

The resolution is that Di Ogum is the
artist's name and spiritual/artistic
connection.  The individual object can
still represent the \emph{ofa} of Ox\'ossi.
The work in iron also gives a separate
material connection with Ogum.

\section{Books encountered at MUNCAB}

The following books were photographed or
examined:

\begin{itemize}
\item Danillo Silva Barata,
\emph{Pedagogia da ancestralidade: f\'e,
educa\c{c}\~ao e lideran\c{c}a no Il\^e Ax\'e
Op\^o Afonj\'a}, 2025,
ISBN 978-85-5971-143-1;
\item Danillo Silva Barata,
\emph{Alberto Pitta: F\'unf\'un D\'ud\'u},
2025, ISBN 978-85-5971-149-3;
\item Danillo Silva Barata,
\emph{J. Cunha e o Carnaval Negro}, 2024,
ISBN 978-65-88622-04-9;
\item Zulu Ara\'ujo,
\emph{Apenas um cidad\~ao}, 2025,
ISBN 978-85-5330-070-9;
\item Tha\'\i s Salves de Brito and others,
\emph{Bemb\'e do Mercado: processo de
instru\c{c}\~ao de registro},
ISBN 978-85-5971-146-2;
\item Mateus Aleluia Lima, organizer,
\emph{Na\c{c}\~oes do Candombl\'e: Jeje},
2025, ISBN 978-85-5971-147-9.
\end{itemize}

These books occupy different places in the
project.  The Afonj\'a book is a study of a
specific Ketu/Nag\^o house.  The Jeje book
is a deliberately particular study of
another Candombl\'e matrix.  The Pitta and
J. Cunha books belong primarily to art,
textile, carnival, and visual culture.  The
Zulu Ara\'ujo book is autobiographical and
belongs to the social and political history
of the Black movement.  The Bemb\'e book
concerns religion, public space, heritage,
and the Rec\^oncavo.

\section{Orun Aiye}

A film cover read
\emph{Orun Aiye---A cria\c{c}\~ao do mundo}.
It named Carlinhos Brown, Carlos Bet\~ao,
Jorge Washington, and Jo\~ao Miguel.

The title called attention to the
opposition and relation between the
spiritual realm, \emph{Orun}, and the
earthly realm, \emph{Aiye}.  It also
created the later need to distinguish
\emph{Orun} from \emph{Olorun}.

Jorge Washington was later identified as a
longstanding actor of the Bando de Teatro
Olodum.  Any assignment of the film's roles
should be checked against its official
credits before entering the permanent
notes.

\section{Nadia Taquary and Iroko}

A Bienal panel described N\'adia Taquary's
\emph{Ir\^ok\^o: A \`arvore c\'osmica}, 2025.
The panel connected the work with feminine
power, African-derived technologies and
narratives, bead strings in the colours of
Iroko, ancestral knowledge, life cycles,
time, and ancestry.

It stated that Iroko is cultivated in
Brazil through the gameleira, marked in
terreiros by a white flag.  It also said
that the Iyami rested on the primordial
tree.

The word \emph{Iyami} called attention.
The panel translated it as
``feiticeiras''.  That translation is a
museum choice that needs conceptual
qualification; it should not simply be
identified with the European category
``witch''.

The comparison with Circe is useful as a
comparative question about female power,
transformation, fear, and translation.  It
is not evidence of a historical relation
between Greek and Yoruba traditions.

\section{Tincoa and music}

The exhibition \emph{Tinco\~a} joined two
referents:

\begin{itemize}
\item the tinco\~a, a bird whose call the
curatorial text associated with memory and
ancestral messages;
\item Os Tinco\~as, the Bahian vocal group.
\end{itemize}

This encounter established music as a
fourth axis of the trip, alongside history,
religion, and language.  The museum text's
claim that Bahia is the most Black state
outside Africa should be treated as a
slogan to verify statistically, not as an
established comparative fact.

\section{Sandro Aiye}

In the sculpture garden, a label identified
Sandro Aiy\^e's \emph{B\'ar\`a II}, 2025,
as a carving in demolition wood with
beads.

The larger exhibition was
\emph{Pad\^e On\~a---Encontrar Caminhos}.
Its official framing mobilizes Exu as
movement, mediation, circulation, and the
opening of paths.  This gives a better
context for the work than identifying it
only from a trident-like attribute.

\section{African map and Akan label}

A decorative mask label read
``Povos Akan---\'Africa Central''.  A map in
the same museum placed the Akan region in
Ghana and C\^ote d'Ivoire, that is, West
Africa.

This internal inconsistency must be
preserved.  The map also distinguished
Yoruba regions near Nigeria and Benin and
several Central African peoples near Congo
and Angola.  It made visible that cultural
regions and modern national borders are
not the same kind of object.

\section{Zumbi and Nzambi}

A monument outside represented Zumbi dos
Palmares, the historical leader associated
with Palmares and Black resistance to
slavery.

The name prompted a memory of Nzambi or
Nzambi a Mpungu, a supreme creator name in
Central African/Kongo traditions.  The two
referents must not be confused.  A possible
linguistic relationship also should not be
asserted merely from resemblance.

\section{Mãe Stella and the manifesto}

The MUNCAB reproduced the 1983 manifesto
\emph{O p\'ublico e o povo do Candombl\'e}.
The photographed panel attributed it to
Stella de Ox\'ossi, Menininha do Gantois,
Tet\^e de Ians\~a, Olga de Alaketu, and
Nicinha do Bogum, and dated the newspaper
publication to August 12, 1983.

The manifesto rejected descriptions of
Candombl\'e as a sect or primitive animism.
It treated syncretism as historically tied
to slavery and submission, defended the
autonomy of Candombl\'e, and criticized
folklorization, commercialization, press
sensationalism, and touristic consumption,
including the Lavagem do Bonfim.

The important nuance is that the manifesto
did not deny that syncretism had existed.
It interpreted its continuation as a
political and religious problem after the
conditions that had made concealment a
strategy of survival.

The M\~ae Stella timeline recorded:

\begin{itemize}
\item birth in Salvador in 1925;
\item initiation by M\~ae Senhora in 1939
and the name Od\'e Kayod\^e;
\item decades of work in nursing and public
health;
\item election as the fifth ialorix\'a of
Op\^o Afonj\'a in 1976;
\item travel to Nigeria in 1981;
\item the public anti-syncretism position
in 1983;
\item international conferences and a
journey to Benin with Pierre Verger;
\item books, public honours, the federal
protection of the terreiro, and cultural
and educational projects;
\item election to the Academia de Letras da
Bahia in 2013;
\item death in 2018, described as a return
to Orun.
\end{itemize}

This made clear why the Op\^o Afonj\'a book
is important.  M\~ae Stella was not simply
a famous ritual leader but a major public
and intellectual actor in modern Brazilian
Candombl\'e.

\section{Catholic room at MUNCAB}

A panel on Marian devotion listed Our Lady
under the invocations of the Conception,
the Rosary, Sorrows, the Good Death, and
Carmel.

This prompted two basic distinctions.
These are different titles of Mary, not
different goddesses.  The Immaculate
Conception concerns Mary's own conception
without original sin; it is not the same as
the virginal conception of Jesus.

The title Our Lady of Carmel comes from the
Carmelite tradition and Mount Carmel.  It
also explains part of Salvador's Carmo
toponymy.

\section{Saint George plant}

A garden label identified
\emph{Dracaena angolensis} by the popular
name ``lan\c{c}a-de-S\~ao-Jorge''.

The name joins a Christian warrior saint,
a lance-shaped plant, and popular ideas of
protection.  Any mapping from Saint George
to Ogum or Ox\'ossi must be specified by
region and religious house rather than
presented as a universal equivalence.

\section{Casa do Benin}

A textile map showed the modern Republic
of Benin as a long, narrow state between
Togo and Nigeria, with a short Atlantic
coast.  K\'etou appeared as an actual place,
not only as a Brazilian religious label.

The museum's central question is the
historical Bahia--Benin relation.  It is
not merely a generic museum ``about
Africa''.  Pierre Verger's image of flux
and reflux gives the guiding problem:
what crossed from Africa to Bahia, what was
transformed in Bahia, and what later moved
back across the Atlantic?

Emo de Medeiros's
\emph{Le vaisseau de la Libert\'e---O navio
da Liberdade}, 2025, showed an invented,
brightly patterned Atlantic ship.  The
work should be read as contemporary
reworking, not mistaken for a photograph
of a historical vessel.  The claim that it
used generative AI was made in conversation
but was not established by the photographed
label and should not enter the notes without
a reliable source.

\section{Mirongar}

A Casa do Benin wall text titled
\emph{Mirongar} described racism as a
historical regime of discourse and image.
It proposed ``mirongar'' images and
statements by working between what can be
shown and what must remain guarded.

The text's own formulation made
\emph{mironga} signify mystery, secret,
strategy, and counter-visual practice.  It
also connected this practice with escape
routes and networks of affection.

The later etymological question is separate:
what African language supplied the word,
and through what semantic history did it
enter Brazilian Portuguese?  The Aulete
entry points through \emph{milonga} to a
Kimbundu form, but this needs specialist
verification before a precise etymology is
fixed.

\section{Questions retained}

The visit generated the following research
questions:

\begin{itemize}
\item How did distinct African peoples and
languages become Candombl\'e nations in
Brazil?
\item What precisely distinguishes Yoruba,
Nag\^o, and Ketu?
\item How do Candombl\'e, Umbanda, and Cuban
Santer\'\i a differ historically?
\item What are the religious and linguistic
relations among Orun, Aiye, Olorun, ori,
and orix\'a?
\item How do beads operate as material,
ornament, religious sign, and consecrated
object?
\item How should one compare Iyami and
Circe without inventing a genealogy?
\item How did Salvador and the Rec\^oncavo
become such important centres of
Afro-Brazilian life?
\item How should museum labels be checked
when they contradict their own maps?
\end{itemize}

\bibliography{my,salvador}
\bibliographystyle{amsalpha}

\end{document}
